% !TEX root = ../main.tex
\chapter{Laboratory Safety Protocols}

\section{Electrical Hazards \& Human Physiology}
Electronics laboratory environments present personnel with hazardous voltages, arc flash hazards, thermal sources, toxic chemical vapors during soldering, and sensitive electrostatic discharge (ESD) risks. Compliance with occupational health and safety standards \cite{yildiz2010approach, bicol_university_safety} is essential to maintain a zero-injury environment.

\subsection{Current vs. Voltage: Physiological Impact}
A primary tenet of electrical safety is that \textbf{electric current}, rather than voltage alone, determines physiological trauma. The human body behaves as an electrolytic volume conductor whose internal impedance varies drastically depending on tissue hydration, contact surface area, and skin moisture.

These physiological thresholds were empirically quantified through the landmark biomedical research of Professor Charles F. Dalziel at the University of California, Berkeley \cite{dalziel1968lethal}. Dalziel's extensive experimental trials on human sensory perception and muscular tetany established the definitive criteria for current tolerance across commercial power frequencies ($50\text{--}60\text{ Hz}$). His findings subsequently served as the foundational basis for international regulatory standards, including the International Electrotechnical Commission Technical Report \textbf{IEC TR 60479-1} \cite{iec60479_1}, \textbf{OSHA Publication 3075} (\textit{Controlling Electrical Hazards}) \cite{osha3075}, and \textbf{IEEE Standard 80} \cite{ieee80}.

\begin{table}[htbp]
    \centering
    \caption{Physiological Effects of 60 Hz Alternating Current on the Human Body}
    \label{tab:shock_thresholds}
    \begin{tabularx}{\textwidth}{p{2.8cm} p{4.2cm} X}
        \toprule
        \textbf{Current Level} & \textbf{Threshold Name} & \textbf{Physiological Reaction} \\
        \midrule
        $1\text{ mA}$ ($0.001\text{ A}$) & Perception Threshold & Faint tingling sensation at contact point; generally harmless. \\
        $10\text{ mA}$ ($0.010\text{ A}$) & ``Let-Go'' Threshold & Involuntary muscle contraction; motor control lost; victim cannot release energized conductor. \\
        $20\text{--}30\text{ mA}$ & Respiratory Paralysis & Severe thoracic muscle tetany; breathing arrests; rapid asphyxiation if loop is sustained. \\
        $\ge 100\text{ mA}$ ($0.1\text{ A}$) & Ventricular Fibrillation & Uncoordinated cardiac twitching; blood circulation ceases; fatal within seconds if not defibrillated. \\
        \midrule
        \multicolumn{3}{l}{\footnotesize\textit{Sources: Dalziel \& Lee (IEEE) \cite{dalziel1968lethal}, IEC TR 60479-1 \cite{iec60479_1}, and OSHA Pub.~3075 \cite{osha3075}.}} \\
        \bottomrule
    \end{tabularx}
\end{table}
\subsection{Body Resistance \& Environmental Variables}
Intact, calloused dry skin offers a resistance between $100\text{ k}\Omega$ and $500\text{ k}\Omega$. However, when skin is broken, punctured, or saturated with perspiration, skin resistance collapses below $1\text{ k}\Omega$. Under wet conditions, voltages as low as \textbf{50 V} can deliver lethal currents exceeding $50\text{ mA}$ across the chest. Furthermore, high-voltage test equipment (such as cathode ray tubes or inductive supplies) can exceed $10{,}000\text{ V}$, making liquid spills an immediate life hazard.

\begin{safetybox}
\textbf{CRITICAL SAFETY RULES FOR THE BENCH:}
\begin{itemize}[leftmargin=1.2em]
    \item \textbf{Remove All Conductive Jewelry:} Metallic rings, bracelets, watches, and chains drastically lower skin contact resistance and can cause severe flash burns or direct shock paths.
    \item \textbf{The ``One-Hand Rule'':} When probing circuits with potentially dangerous potentials, keep your non-dominant hand in a pocket or behind your back. This prevents forming a closed circuit loop across both arms through the thoracic cavity and heart.
    \item \textbf{No Food or Beverages:} Drinks near workbenches are strictly prohibited to prevent conductive fluid spills into energized test instruments.
    \item \textbf{Verify Zero Potential:} Treat all circuits as energized until de-energizing and physically verifying zero potential with a calibrated Digital Multimeter (DMM).
\end{itemize}
\end{safetybox}

\section{Facility Safety, Grounding, \& Fire Protection}
Workstations must utilize three-prong grounded power distribution equipped with Ground Fault Circuit Interrupters (GFCIs) or isolation transformers for AC power line equipment. Workstations must maintain clear access to:
\begin{itemize}[leftmargin=1.5em]
    \item Emergency master power disconnect shutoffs.
    \item Class C or Class ABC dry chemical/CO$_2$ fire extinguishers rated for energized electrical equipment. Never apply water to an electrical fire.
    \item Standard first aid stations and eye-wash facilities.
\end{itemize}

\section{Electrostatic Discharge (ESD) Prevention}
Friction from clothing and walking across polymer flooring can generate static charges exceeding \textbf{5,000 V}. Human sensory perception detects static discharges only above 3,500--4,000 V, yet fragile semiconductor junction gates (especially MOSFETs) can suffer dielectric breakdown from pulses as low as 20--50 V.

Technicians must wear grounded anti-static wrist straps incorporating a minimum $500\text{ k}\Omega$ (standard $1\text{ M}\Omega$) safety series resistor. This resistor safely drains electrostatic accumulation to earth ground while providing adequate impedance to limit current to safe sub-milliampere levels in case the operator touches an energized mains conductor.

\section{Workbench Housekeeping \& Shutdown Protocol}
At the end of every laboratory session, technicians must transition the bench into ``Secure Mode'':
\begin{enumerate}[leftmargin=1.5em]
    \item Power down all signal generators, oscilloscopes, and DC power supplies.
    \item Reset variable voltage and current knobs to zero.
    \item Clean and tin soldering iron tips, confirm irons are seated in insulated stands, and turn off iron mains power.
    \item Return passive components to sorted bin cabinets and store leads and tools in labeled drawers.
\end{enumerate}

\section{Laboratory Power Tool Operation \& Mechanical Safety}
Power tools commonly utilized in electronics laboratories---such as benchtop PCB drill presses, rotary cutting tools, and hot air rework stations---introduce severe mechanical and thermal hazards alongside electrical risks. Safe operation requires strict adherence to mechanical protocols:

\begin{itemize}[leftmargin=1.5em]
    \item \textbf{Workpiece Clamping:} Never hold a printed circuit board or chassis enclosure by hand during drilling or cutting operations. The rotational torque of high-speed bits can snatch the workpiece, causing violent projectile ejection or deep lacerations. Workpieces must always be immobilized using a mechanical drill press vise or PCB fixture.
    \item \textbf{Entanglement Prevention \& Dress Code:} Rotating spindles present an immediate entanglement hazard. Personnel must strictly observe a \textbf{zero-gloves rule} when operating drills and rotary tools, as fabric can become ensnared in spinning chucks. Long hair must be tied back securely, long sleeves rolled above the elbows, and lanyards, loose jewelry, and neckties removed prior to operating machinery.
    \item \textbf{Particulate and Eye Protection:} Machining copper-clad FR4 (flame-retardant fiberglass epoxy) generates microscopic glass dust and sharp copper burrs. Operators must wear ANSI Z87.1-rated safety glasses with side shields and an N95 particulate respirator to prevent corneal abrasion and inhalation of toxic fiberglass particulates.
    \item \textbf{Tool Maintenance \& Disconnection:} Always disconnect power from rotary tools and drill presses before changing drill bits, cutting discs, or sanding drums. Never leave an energized tool unattended on a workbench.
\end{itemize}

\section{PCB Etchant Chemical Safety: Storage, Neutralization \& Disposal}
Wet chemical etching processes commonly employ acidic oxidizers such as Ferric Chloride ($\text{FeCl}_3$) or Ammonium/Sodium Persulfate to dissolve unwanted copper laminate. Because these solutions are highly corrosive and toxic, handling them demands strict compliance with occupational safety and environmental regulations:

\subsection{Storage Protocols}
\begin{itemize}[leftmargin=1.5em]
    \item \textbf{Chemical Compatibility:} Ferric chloride and persulfate solutions must be stored strictly in acid-resistant, heavy-duty High-Density Polyethylene (HDPE), polypropylene, or fluoropolymer containers with hermetic seals. \textbf{Never store or mix etchants in metallic containers or sinks}, as the acid reacts vigorously with metals, generating corrosive fumes and risking container rupture.
    \item \textbf{Environment and Segregation:} Store etchant bottles inside a dedicated, ventilated acid/corrosive chemical cabinet. Keep containers isolated from organic solvents, flammable materials, strong bases (which cause exothermic reactions), and direct ultraviolet sunlight.
    \item \textbf{Container Labeling:} Every container must bear GHS-compliant hazardous labeling clearly stating the chemical name, concentration, preparation date, and specific hazard warnings (Corrosive, Environmental Hazard).
\end{itemize}

\subsection{Neutralization and Disposal Protocols}
\begin{itemize}[leftmargin=1.5em]
    \item \textbf{Strict Prohibition of Drain Disposal:} \textbf{Spent etchant solutions must NEVER be poured into domestic, commercial, or municipal plumbing.} Dissolved cupric ions ($\text{Cu}^{2+}$) are highly biocidal, permanently disrupt biological wastewater treatment facilities, and rapidly dissolve copper and steel drainage piping.
    \item \textbf{Pre-Disposal Neutralization:} Before transfer, spent acidic etchant must be slowly neutralized by gradually adding a weak base such as Sodium Bicarbonate (baking soda, $\text{NaHCO}_3$) or Calcium Carbonate ($\text{CaCO}_3$) in a well-ventilated fume hood until carbon dioxide effervescence ceases ($\text{pH} \approx 7.0\text{--}8.0$). This reaction converts toxic soluble copper ions into insoluble copper carbonate/hydroxide precipitate.
    \item \textbf{Collection and Institutional Hand-off:} Allow the precipitated copper sludge to settle. Decant the neutral liquid and collect the sludge in designated hazardous chemical waste carboys labeled ``Spent Copper Etchant Sludge''. Submit containers to the University Environmental Safety Office or a licensed hazardous chemical disposal contractor in full compliance with local environmental legislation (e.g., DENR standards).
\end{itemize}
